\section{Introduction}
\label{sec:introduction}

% src: 47 CFR 9.10(i)(2)(i) (dispatchable location or x/y within 50 m for 80 % of wireless 911 calls within 6 years; law.cornell.edu, verified 2026-10-30); Italiano et al. 2025 p. 5 (emergency-call requirements drove cellular localization; "support emergency services and foster location-based applications"; asset tracking), p. 2 (5G "complementing GNSS ... in outdoor positioning and navigation"); Groves 2011 abstract ("GPS is unreliable in dense urban areas ... buildings block the signals from many of the satellites")
Locating a user equipment~(UE) from the cellular network is a regulatory requirement and an enabler of location-based services. In the United States, wireless carriers must locate \SI{80}{\percent} of wireless 911 calls to within \SI{50}{\meter}~\cite{cfr47_9_10}, and cellular positioning also serves applications such as asset tracking and navigation~\cite{italiano2025tutorial}. The need is sharpest where satellite positioning is weakest, as in dense urban areas, where buildings block the signals of many satellites~\cite{groves2011shadow}.

% src: 3GPP TS 38.215 §5.1 (UE: SS-RSRP, SS-SINR), §5.2 (gNB: UL AoA, TADV); TS 38.305 §4.3.10 (NR E-CID: "the UE generally is not expected to make additional measurements for the sole purpose of positioning"), §8.9.1 (E-CID inputs); verified on itecspec.com 2026-10-30
The cellular network is well placed to meet this need, because it already observes every connected UE. In 5G New Radio~(NR), the UE measures and reports its signal-to-interference-plus-noise ratio~(SINR) and its received signal strength~(RSS), reported as the reference signal received power~(RSRP), while the base station~(BS) measures the uplink angle of arrival~(AoA) and the timing advance~\cite{3gpp38215}. The standardized enhanced cell-ID method uses these measurements for positioning without requiring any made for that purpose alone~\cite{3gpp38305}, so the network can position a UE without an added receiver, sensor or infrastructure.

% src: Italiano et al. 2025 p. 14 (RSS path-loss model Eqs. (14)-(15); power-based measurements are "reasonable distance indicators only if the BS is near to the UE", shadowing and multipath), p. 22 ("DL-TDOA requires precise synchronization among the BSs"; RTT does not)
The difficulty is accuracy. RSS depends mainly on the distance to the BS and is perturbed by shadowing and multipath, so it is a reasonable distance indicator only near the BS~\cite{italiano2025tutorial}, and it carries no information on the bearing. Geometric positioning recovers the bearing by combining time or angle measurements from several BSs, and its time-difference methods further require precise synchronization among them~\cite{italiano2025tutorial}. A single serving BS, however, provides at most one range and, where an angle is available, one bearing; without a reliable angle, a single report leaves the position under-determined. This case would benefit most from a solution that adds no infrastructure. We address it by using the history of a UE's reports, not only the current one.

\subsection{The Channel History Hypothesis}
\label{subsec:ambiguity-hypothesis}

% src: FINAL_REPORT.md §1.1, §1.2, §3.2
The hypothesis evaluated in this work is that a moving UE's sequence of reports carries more positional information than any single report, and that this holds regardless of the estimator applied on top of it. Because the radio channel is spatially consistent, consecutive reports of a moving UE are correlated samples of the channel along its path, and a window of them shows how the measurements change as the UE moves, which no single report contains.

The ambiguity of a single report is one reason to expect such a gain. RSS depends mainly on the distance to the BS, so a single reading constrains the range, up to shadow fading, but not the bearing: positions on a common range ring around the BS are indistinguishable to a single-snapshot estimator, which we refer to as the \emph{distance-ring ambiguity}. SINR and AoA can each supply bearing information, but both are conditional, since SINR depends on the placement of interfering BSs and AoA on UE hardware that legacy single-antenna devices lack. Two positions on the same ring, however, produce different sequences of reports, because the channel around each evolves differently as the UE moves. A history of reports may therefore resolve what a single report cannot, without additional infrastructure; this work tests whether it does so across propagation settings, BS placements and estimators.

\subsection{Contributions}
\label{subsec:contributions}

% src: CLAIM_EVIDENCE_MAP.md (2026-10-24) core message and C1-C9; studies G, R, M as defined there
We evaluate this hypothesis in three simulated studies of a single 5G NR BS: a controlled $15\times15$ grid on which the exact cell is classified, $25\times25$ regression runs that pool five device profiles, and a macro cell in which every model is tested on users not seen in training. Using the history of a user's reports, rather than the current report alone, consistently lowers the error of the best model across propagation settings, BS placements and model families. The gain depends on the position information a single report carries and on whether that information is shared between users: where one report carries nothing that transfers to new users, as for single-antenna devices in the uniform-NLOS macro cell, whose users are each simulated as a separate channel realization, history adds nothing, whereas in a simulated environment shared by all users, history lowers the error of unseen users, under uniform NLOS as well.

Specifically, the contributions of this work are the following: (i) a consistent history gain, measured for the best model in every propagation setting and BS placement of the three studies and in all seven model families of the macro cell, and robust to propagation, blockage and interference, thereby making the input, rather than the estimator, the deciding design choice; (ii) evidence that sequence models exploit long windows where tree ensembles given the same window cannot, together with a test showing that the order of the reports carries part of that gain on the grid; (iii) an account of the macro-cell gain as the averaging of noisy angle reports, and of how it shrinks when the angle error persists across reports; and (iv) a delimitation of where history adds nothing, namely where a single report carries no position information that transfers between users, together with a test showing that this depends largely on the simulation: when all users share one simulated environment, history lowers the error of unseen users, also under uniform NLOS, and differences between consecutive reports become informative.

\noindent\textbf{Organization.} Section~\ref{sec:related-work} reviews related work. Section~\ref{sec:system-model} formulates the problem and what history can add to a single report. Section~\ref{sec:setup} describes the simulated deployments, devices, features, estimators and evaluation protocol. Section~\ref{sec:results-gain} presents the history gain across the three studies and its robustness to propagation, blockage and interference; Section~\ref{sec:results-window} then compares how tree ensembles, sequence models and a tracking filter use the window, and Section~\ref{sec:results-mechanism} identifies what history reads in the macro cell. Section~\ref{sec:results-no-gain} examines the one setting in which history adds nothing. Section~\ref{sec:discussion} discusses the results and their limitations; Section~\ref{sec:conclusion} concludes.


\subsection{Related Work}
\label{sec:related-work}

% src: FINAL_REPORT.md §2, "Fingerprinting" paragraph
Classical RSS fingerprinting matches a measurement snapshot against a surveyed database~\cite{yang2013_rssitocsi}; single-site multipath fingerprinting establishes that one BS already observes enough structure to localize without triangulation~\cite{kupershtein2013}, the principle this work extends to a moving user. Fingerprinting's characteristic failure mode is poor discrimination between points at equal range from the BS, which is exactly the distance-ring ambiguity a history of reports targets.
The static $k$-nearest-neighbor~($k$-NN) regressor at $h=0$ represents this family throughout this work. % src: FINAL_REPORT.md §2 (k-NN h=0 baseline)

% src: FINAL_REPORT.md §2, "Geometric methods" paragraph
Geometric methods recover position from an explicit system built on time-of-arrival, time-difference-of-arrival, or AoA estimates~\cite{italiano2025tutorial}.
These methods require either multiple synchronized BSs or a high-quality angle estimate, so the single-serving-BS, single-angle-source deployment studied here is precisely the regime in which the geometric system is underdetermined: for a device that reports no angle at all, the available geometry constrains range but leaves bearing unresolved anywhere along the range ring.

% src: FINAL_REPORT.md §2, "Sequence models for CSI" paragraph
Recurrent and dynamic-system formulations have likewise been applied to channel state information~(CSI)-based tracking~\cite{fang2010dynamic,yu2020rnn,shi2018tracking}, and a related line shows that the temporal autocorrelation of RSS improves fingerprinting accuracy when the signal is treated as a stationary process observed at a fixed location~\cite{wang2016temporal,tian2018temporal}. That signal differs from the one used here: it is order-invariant and aggregated across access points, so it survives for a stationary receiver and vanishes for a moving one, whereas the channel-history window used here is ordered, is read from a single BS, and encodes exactly the motion the in-place statistic discards. These works propose a tracking architecture and evaluate it as a whole, so the benefit of the temporal input cannot be separated from that of the architecture. In this work, we instead hold each estimator fixed and vary only the history window it receives, thereby isolating the contribution of history, and repeat the comparison across model families.
